\chapter*{Preface}
Toss a fair coin a million times. The law of large numbers says the proportion of heads is close to one half, and the central
limit theorem says how close: within about one thousandth, with the bell curve describing the error. These are the first two
answers of probability, and they raise every question of this book. How good is the bell-curve approximation, and what
replaces it when the variables have no variance? How unlikely is it to see sixty percent heads, exactly, on the exponential
scale? Why are so many functions of many independent variables almost constant? What does the whole path of the walk look like
when you step back far enough to see it as a curve? And what can be said when the variables are not independent at all, but
come from a single deterministic system observed over time?

The first part of the book proves the limit theorems in their sharp forms. Zero--one laws and random series open the way;
then the strong law with Etemadi's proof and Kolmogorov's three-series theorem; weak convergence in spaces of functions, with
Donsker's invariance principle, which turns every random walk with finite variance into Brownian motion, and the law of the
iterated logarithm; the central limit theorem with the Lindeberg--Feller converse and the Berry--Esseen bound; and the
infinitely divisible and stable laws, which are all the possible limits of sums.

The second part measures. Stein's method gives explicit distances to the normal and Poisson laws with one elegant identity.
Cram\'er's and Sanov's theorems compute the probabilities of rare events on the exponential scale. Concentration of measure,
from McDiarmid's inequality to Talagrand's, explains why the longest increasing subsequence of a random permutation and the
chromatic number of a random graph hardly fluctuate, and why a thousand points in a million dimensions can be projected to a
few hundred without distorting their distances.

The third part leaves independence behind. Martingales return in full strength, with de Finetti's theorem on exchangeable
sequences and Kakutani's dichotomy. Ergodic theory replaces independence by invariance: Birkhoff's theorem computes the
statistics of the continued fraction of almost every real number, and Furstenberg's method proves that every set of integers
of positive density contains three terms in arithmetic progression. The book ends where probability meets arithmetic:
mod-$\varphi$ convergence, a sharper form of the central limit theorem that sees the number of prime factors of an integer and
the number of cycles of a permutation in the same light, and the Poisson--Dirichlet law, which describes both the cycle
lengths of a large random permutation and the sizes of the prime factors of a large random integer.

The aim throughout is to understand why the theorems are true, and to see them at work. Every result is proved, with one exception: the Selberg--Delange theorem
of Chapter~11 belongs to analytic number theory, and is quoted there for comparison and proved in Volume~14. Pause and think
boxes ask you to guess before reading on. Each chapter ends with exercises, graded from warm-ups ($\star$) to small gems
($\star\star\star$), each with a hint and a complete solution that first explains how to find the idea and then writes it up.

The book assumes the probability of Volume~3, the measure and integration of Volume~6 and the Hilbert spaces of Volume~7. Its
beauty is in the universality of its answers: the same few laws, Gaussian, Poisson, stable, Brownian, Poisson--Dirichlet,
appear whatever the details, and the chapters that follow try to show why.
